\subsection{Hyperparameter Tuning}

The hyperparameter tuning phase is designed to assess the \emph{stability} and
\emph{robustness} of the proposed two-stage formulation, rather than to maximize
performance through exhaustive search.

Unlike end-to-end models, the Two-Stage classifier is governed by a small and
interpretable set of hyperparameters controlling the separation between
near-deterministic and ambiguous instances.

\paragraph{Tuned Parameters}
The only hyperparameters affecting the Two-Stage model are:
\begin{itemize}
	\item the minimum rule purity threshold $\tau_p$,
	\item the minimum rule support threshold $\tau_s$,
	\item the regularization strength $C$ of the Stage~2 linear classifier.
\end{itemize}

\paragraph{Controlled Grid Analysis}
A controlled grid search is first conducted on the development split to explicitly
analyze the trade-off between rule coverage and overall macro-F1 performance.
The explored ranges are:
\[
\tau_p \in \{0.90, 0.95\}, \quad
\tau_s \in \{20, 30, 40\}, \quad
C \in \{0.5, 1.0, 1.5\}.
\]

Across all configurations, Stage~1 resolves approximately $16\%$--$20\%$ of the
dataset with near-perfect purity, while Stage~2 handles the remaining ambiguous
documents. Performance gains saturate once the most deterministic patterns are
captured, indicating that improvements are driven by structural decomposition
rather than fine-grained parameter tuning.

\paragraph{Coverage--Performance Trade-off}
Figure~\ref{fig:coverage_f1_tradeoff} illustrates the relationship between the
fraction of samples resolved by Stage~1 and the resulting macro-F1 score.
The plot highlights a stable operating region where increased coverage yields
consistent gains without degrading generalization.

\begin{figure}[t]
	\centering
	% INSERT IMAGE HERE:
	% Scatter plot with:
	% x-axis = Rule coverage (% of samples resolved by Stage 1)
	% y-axis = Macro-F1 score
	% points = different (tau_p, tau_s) configurations
	\includegraphics[width=0.85\linewidth]{coverage_vs_macroF1.pdf}
	\caption{Coverage--performance trade-off for the Two-Stage model.
	Each point corresponds to a different $(\tau_p, \tau_s)$ configuration.}
	\label{fig:coverage_f1_tradeoff}
\end{figure}

\paragraph{Best Configurations}
Table~\ref{tab:two_stage_tuning} reports the top-performing configurations.
All high-performing settings cluster in a narrow region of the parameter space,
confirming the robustness of the Two-Stage formulation.

\begin{table}[t]
	\centering
	\caption{Top Two-Stage configurations on the development split.}
	\label{tab:two_stage_tuning}
	\begin{tabular}{cccc}
		\toprule
		$\tau_p$ & $\tau_s$ & Coverage (\%) & Macro-F1 \\
		\midrule
		0.95 & 30 & $\sim$16 & 0.721 \\
		0.95 & 40 & $\sim$15 & 0.721 \\
		0.90 & 30 & $\sim$19 & 0.720 \\
		\bottomrule
	\end{tabular}
\end{table}

\paragraph{Bayesian Optimization Verification}
A lightweight Bayesian optimization (Optuna) is subsequently applied over the
same parameter space to verify that no alternative high-performing regions are
missed. The optimizer consistently converges to the same configurations identified
by the controlled grid search, confirming that the Two-Stage model does not rely
on fragile or dataset-specific tuning.

Overall, hyperparameter tuning confirms that the performance gains of the
Two-Stage approach stem from an explicit separation between deterministic and
ambiguous regimes, rather than from aggressive parameter optimization.
